% Options for packages loaded elsewhere
% Options for packages loaded elsewhere
\PassOptionsToPackage{unicode}{hyperref}
\PassOptionsToPackage{hyphens}{url}
%
\documentclass[
  ignorenonframetext,
  aspectratio=169,
  russian,
]{beamer}
\newif\ifbibliography
\usepackage{pgfpages}
\setbeamertemplate{caption}[numbered]
\setbeamertemplate{caption label separator}{: }
\setbeamercolor{caption name}{fg=normal text.fg}
\beamertemplatenavigationsymbolsempty
% Prevent slide breaks in the middle of a paragraph
\widowpenalties 1 10000
\raggedbottom
\AtBeginPart{
  \frame{\partpage}
}
\AtBeginSection{
  \ifbibliography
  \else
    \frame{\sectionpage}
  \fi
}
\AtBeginSubsection{
  \frame{\subsectionpage}
}
\usepackage{iftex}
\ifPDFTeX
  \usepackage[T2A]{fontenc}
  \usepackage[utf8]{inputenc}
  \usepackage{textcomp} % provide euro and other symbols
\else % if luatex or xetex
  \usepackage{unicode-math} % this also loads fontspec
  \defaultfontfeatures{Scale=MatchLowercase}
  \defaultfontfeatures[\rmfamily]{Ligatures=TeX,Scale=1}
\fi
\usepackage{lmodern}

\usetheme[progressbar=frametitle,sectionpage=progressbar,numbering=fraction]{metropolis}
\usefonttheme{serif} % use mainfont rather than sansfont for slide text
\ifPDFTeX\else
  % xetex/luatex font selection
  \setmainfont[Ligatures=TeX]{PT Serif}
  \setsansfont[Ligatures=TeX,Scale=MatchLowercase]{PT Sans}
  \setmonofont[Scale=MatchLowercase,Scale=0.9]{PT Mono}
\fi
% Use upquote if available, for straight quotes in verbatim environments
\IfFileExists{upquote.sty}{\usepackage{upquote}}{}
\IfFileExists{microtype.sty}{% use microtype if available
  \usepackage[]{microtype}
  \UseMicrotypeSet[protrusion]{basicmath} % disable protrusion for tt fonts
}{}
\makeatletter
\@ifundefined{KOMAClassName}{% if non-KOMA class
  \IfFileExists{parskip.sty}{%
    \usepackage{parskip}
  }{% else
    \setlength{\parindent}{0pt}
    \setlength{\parskip}{6pt plus 2pt minus 1pt}}
}{% if KOMA class
  \KOMAoptions{parskip=half}}
\makeatother


\usepackage{longtable,booktabs,array}
\usepackage{calc} % for calculating minipage widths
\usepackage{caption}
% Make caption package work with longtable
\makeatletter
\def\fnum@table{\tablename~\thetable}
\makeatother
\usepackage{graphicx}
\makeatletter
\newsavebox\pandoc@box
\newcommand*\pandocbounded[1]{% scales image to fit in text height/width
  \sbox\pandoc@box{#1}%
  \Gscale@div\@tempa{\textheight}{\dimexpr\ht\pandoc@box+\dp\pandoc@box\relax}%
  \Gscale@div\@tempb{\linewidth}{\wd\pandoc@box}%
  \ifdim\@tempb\p@<\@tempa\p@\let\@tempa\@tempb\fi% select the smaller of both
  \ifdim\@tempa\p@<\p@\scalebox{\@tempa}{\usebox\pandoc@box}%
  \else\usebox{\pandoc@box}%
  \fi%
}
% Set default figure placement to htbp
\def\fps@figure{htbp}
\makeatother



\ifLuaTeX
\usepackage[bidi=basic,provide=*]{babel}
\else
\usepackage[bidi=default,provide=*]{babel}
\fi
\ifPDFTeX
\else
\babelfont{rm}[Ligatures=TeX]{PT Serif}
\fi
% get rid of language-specific shorthands (see #6817):
\let\LanguageShortHands\languageshorthands
\def\languageshorthands#1{}


\setlength{\emergencystretch}{3em} % prevent overfull lines

\providecommand{\tightlist}{%
  \setlength{\itemsep}{0pt}\setlength{\parskip}{0pt}}



 


\makeatletter
\@ifpackageloaded{caption}{}{\usepackage{caption}}
\AtBeginDocument{%
\ifdefined\contentsname
  \renewcommand*\contentsname{Содержание}
\else
  \newcommand\contentsname{Содержание}
\fi
\ifdefined\listfigurename
  \renewcommand*\listfigurename{Список иллюстраций}
\else
  \newcommand\listfigurename{Список иллюстраций}
\fi
\ifdefined\listtablename
  \renewcommand*\listtablename{Список таблиц}
\else
  \newcommand\listtablename{Список таблиц}
\fi
\ifdefined\figurename
  \renewcommand*\figurename{Рисунок}
\else
  \newcommand\figurename{Рисунок}
\fi
\ifdefined\tablename
  \renewcommand*\tablename{Таблица}
\else
  \newcommand\tablename{Таблица}
\fi
}
\@ifpackageloaded{float}{}{\usepackage{float}}
\floatstyle{ruled}
\@ifundefined{c@chapter}{\newfloat{codelisting}{h}{lop}}{\newfloat{codelisting}{h}{lop}[chapter]}
\floatname{codelisting}{Список}
\newcommand*\listoflistings{\listof{codelisting}{Листинги}}
\makeatother
\makeatletter
\makeatother
\makeatletter
\@ifpackageloaded{caption}{}{\usepackage{caption}}
\@ifpackageloaded{subcaption}{}{\usepackage{subcaption}}
\makeatother

\usepackage{bookmark}
\IfFileExists{xurl.sty}{\usepackage{xurl}}{} % add URL line breaks if available
\urlstyle{same}
\hypersetup{
  pdftitle={Доклад},
  pdfauthor={Калашникова О. С., НПИбд-01-23},
  pdflang={ru-RU},
  hidelinks,
  pdfcreator={LaTeX via pandoc}}


\title{Доклад}
\subtitle{Зависимость от пути}
\author{Калашникова О. С., НПИбд-01-23}
\date{2026-04-17}

\begin{document}
\frame{\titlepage}

\renewcommand*\contentsname{Содержание}
\begin{frame}[allowframebreaks]
  \frametitle{Содержание}
  \setcounter{tocdepth}{2}
  \tableofcontents
\end{frame}
\setcounter{tocdepth}{2}
\tableofcontents
}

\section{1. Вводная
часть}\label{ux432ux432ux43eux434ux43dux430ux44f-ux447ux430ux441ux442ux44c}

\begin{frame}{1.1 Вводная часть}
\phantomsection\label{ux432ux432ux43eux434ux43dux430ux44f-ux447ux430ux441ux442ux44c-1}
\begin{columns}[T]
\begin{column}{0.6\linewidth}
\textbf{Актуальность и проблема:}

\begin{itemize}
\tightlist
\item
  Классика: будущее зависит только от настоящего (марковское свойство)
\item
  Реальность: небольшое случайное событие может закрепиться и определить
  исход
\item
  Пример: клавиатура QWERTY (хуже, но победила из-за случайности)
\end{itemize}

\textbf{Главная опасность:} классические эргодические методы дают
неверные прогнозы для таких систем
\end{column}

\begin{column}{0.3\linewidth}
\textbf{Объект исследования:}

Модифицированная урновая модель с зависимостью от всей предыдущей
истории
\end{column}
\end{columns}
\end{frame}

\begin{frame}{1.2 Вводная часть}
\phantomsection\label{ux432ux432ux43eux434ux43dux430ux44f-ux447ux430ux441ux442ux44c-2}
\begin{columns}[T]
\begin{column}{0.3\linewidth}
\textbf{Цель:}

Раскрыть понятие path dependence и его математическую формализацию через
урновую модель
\end{column}

\begin{column}{0.4\linewidth}
\textbf{Задачи:}

\begin{itemize}
\tightlist
\item
  Описать модель и параметры
\item
  Разобрать частные случаи (α = 0, малое α, большое α)
\item
  Показать применение (фильтрация сетей)
\item
  Сравнить с альтернативами
\end{itemize}
\end{column}

\begin{column}{0.2\linewidth}
\textbf{Методы:}

Математическое моделирование, теория вероятностей
\end{column}
\end{columns}
\end{frame}

\begin{frame}{1.3 Введение --- пример QWERTY}
\phantomsection\label{ux432ux432ux435ux434ux435ux43dux438ux435-ux43fux440ux438ux43cux435ux440-qwerty}
\begin{columns}[T]
\begin{column}{0.6\linewidth}
\textbf{Клавиатура QWERTY (1870-е):} - Создана для ЗАМЕДЛЕНИЯ печати
(чтобы не залипали рычажки) - Появилась более удобная раскладка Dvorak

\textbf{Почему QWERTY всё ещё с нами?} - Миллионы людей уже выучили -
Производители выпускали только QWERTY - Рыночные преимущества у
доминирующей раскладки

\textbf{Исход определила историческая случайность, а не качество!}
\end{column}

\begin{column}{0.3\linewidth}
\includegraphics[width=1\linewidth,height=\textheight,keepaspectratio]{image/qwerty.jpg}
\end{column}
\end{columns}
\end{frame}

\section{2. Базовая
модель}\label{ux431ux430ux437ux43eux432ux430ux44f-ux43cux43eux434ux435ux43bux44c}

\begin{frame}{2.1 Урновая модель --- основные понятия}
\phantomsection\label{ux443ux440ux43dux43eux432ux430ux44f-ux43cux43eux434ux435ux43bux44c-ux43eux441ux43dux43eux432ux43dux44bux435-ux43fux43eux43dux44fux442ux438ux44f}
\begin{columns}[T]
\begin{column}{0.5\linewidth}
\textbf{Что у нас есть:} - n = 0, 1, 2, \ldots{} (дискретное время) - a
белых, b чёрных шаров (начально) - Xk = 1 (белый) или 0 (чёрный) - Hn =
(X1, \ldots, Xn) --- вся история - Sn = сумма белых за всю историю
\end{column}

\begin{column}{0.5\linewidth}
\textbf{Правила (классика Пойя):} 1. Вынуть шар пропорционально
количеству 2. Вернуть обратно 3. Добавить c шаров ТОГО ЖЕ цвета

\textbf{Классическая формула:} P = (a + c·Sn) / (a + b + c·n)
\end{column}
\end{columns}
\end{frame}

\begin{frame}{2.2 Что меняем?}
\phantomsection\label{ux447ux442ux43e-ux43cux435ux43dux44fux435ux43c}
\begin{columns}[T]
\begin{column}{1\linewidth}
\textbf{В классике:} вероятность зависит только от Sn (сколько белых
всего), порядок не важен

\textbf{Нам нужно:} чтобы порядок извлечений имел значение!

\textbf{Решение:} добавляем в формулу явный член памяти
\end{column}
\end{columns}
\end{frame}

\section{3. Формула и
параметры}\label{ux444ux43eux440ux43cux443ux43bux430-ux438-ux43fux430ux440ux430ux43cux435ux442ux440ux44b}

\begin{frame}{3.1 Главная формула}
\phantomsection\label{ux433ux43bux430ux432ux43dux430ux44f-ux444ux43eux440ux43cux443ux43bux430}
\begin{columns}[T]
\begin{column}{1\linewidth}
{[} P(X\_\{n+1\}=1 \mid H\_n) =
\frac{W_n + \alpha \cdot S_n + \delta}{W_n + B_n + \beta \cdot n + \gamma}
{]}

\textbf{Упрощённая форма (c=1, β=1, δ=γ=0):}

{[} P = \frac{a + (1+\alpha)S_n}{a+b+n+\alpha S_n} {]}
\end{column}
\end{columns}
\end{frame}

\begin{frame}{3.2 Параметры модели}
\phantomsection\label{ux43fux430ux440ux430ux43cux435ux442ux440ux44b-ux43cux43eux434ux435ux43bux438}
\begin{columns}[T]
\begin{column}{1\linewidth}
\begin{longtable}[]{@{}lll@{}}
\toprule\noalign{}
Параметр & Что делает & Ключевое \\
\midrule\noalign{}
\endhead
\textbf{α ≥ 0} & Сила памяти (главный!) & α=0 → классика, α\textgreater0
→ память \\
β ≥ 0 & Как быстро забываем & β\textgreater1 → память затухает \\
δ, γ & Сглаживание (для малых n) & Убирают нулевые вероятности \\
c & Масштаб добавления & Обычно c=1 \\
\bottomrule\noalign{}
\end{longtable}

\textbf{α --- наш главный инструмент управления зависимостью от пути}
\end{column}
\end{columns}
\end{frame}

\section{4. Частные
случаи}\label{ux447ux430ux441ux442ux43dux44bux435-ux441ux43bux443ux447ux430ux438}

\begin{frame}{4.1 α = 0 --- классическая модель Пойя}
\phantomsection\label{ux3b1-0-ux43aux43bux430ux441ux441ux438ux447ux435ux441ux43aux430ux44f-ux43cux43eux434ux435ux43bux44c-ux43fux43eux439ux44f}
\begin{columns}[T]
\begin{column}{1\linewidth}
{[} P = \frac{a + S_n}{a+b+n} {]}

\begin{itemize}
\tightlist
\item
  Вероятность зависит только от счёта Sn (порядок НЕ важен)
\item
  Mn = (a+Sn)/(a+b+n) --- мартингал
\item
  Предел: Beta(a,b)
\item
  Нет необратимой фиксации
\end{itemize}
\end{column}
\end{columns}
\end{frame}

\begin{frame}{4.2 Малое α \textgreater{} 0 --- «мягкая» память}
\phantomsection\label{ux43cux430ux43bux43eux435-ux3b1-0-ux43cux44fux433ux43aux430ux44f-ux43fux430ux43cux44fux442ux44c}
\begin{columns}[T]
\begin{column}{1\linewidth}
\textbf{Эффекты:} - Начальное поведение --- почти как у классики - Со
временем дисперсия растёт - Вероятность уйти к 0 или 1 повышается (но
медленно)

\textbf{Пример:} α = 0.1, n=100--1000
\end{column}
\end{columns}
\end{frame}

\begin{frame}{4.3 Большое α --- жёсткая фиксация}
\phantomsection\label{ux431ux43eux43bux44cux448ux43eux435-ux3b1-ux436ux451ux441ux442ux43aux430ux44f-ux444ux438ux43aux441ux430ux446ux438ux44f}
\begin{columns}[T]
\begin{column}{1\linewidth}
\textbf{Эффекты:} - α·Sn доминирует уже при небольших n - Начальные
флуктуации катастрофически усиливаются - Высокая вероятность фиксации: -
M\_n → 1 (все белые) - M\_n → 0 (все чёрные)

\textbf{Пример:} α = 10, первые 5--10 шагов определяют всё

\textbf{Это и есть механизм QWERTY!}
\end{column}
\end{columns}
\end{frame}

\section{5. Применение: фильтрация
сетей}\label{ux43fux440ux438ux43cux435ux43dux435ux43dux438ux435-ux444ux438ux43bux44cux442ux440ux430ux446ux438ux44f-ux441ux435ux442ux435ux439}

\begin{frame}{5.1 Задача и урновая аналогия}
\phantomsection\label{ux437ux430ux434ux430ux447ux430-ux438-ux443ux440ux43dux43eux432ux430ux44f-ux430ux43dux430ux43bux43eux433ux438ux44f}
\begin{columns}[T]
\begin{column}{0.5\linewidth}
\textbf{Задача:} - Есть узел сети (степень k, сила s, вес связи w) -
Значима ли связь или случайна?
\end{column}

\begin{column}{0.5\linewidth}
\textbf{Урновая аналогия:} - 1 красный шар (анализируемая связь) - k-1
чёрных (остальные связи) - s извлечений - Модель Пойя →
бета-биномиальное распределение
\end{column}
\end{columns}
\end{frame}

\begin{frame}{5.2 Связь с path dependence}
\phantomsection\label{ux441ux432ux44fux437ux44c-ux441-path-dependence}
\begin{columns}[T]
\begin{column}{1\linewidth}
\textbf{Три ключевых свойства:}

\begin{enumerate}
\tightlist
\item
  \textbf{Положительная обратная связь} --- вынули красный → красных
  стало больше
\item
  \textbf{Немарковость} --- один и тот же состав может быть достигнут
  разными путями
\item
  \textbf{Фиксация} --- случайная связь может стать доминирующей
\end{enumerate}

\textbf{Параметр a в этой модели = аналог α}
\end{column}
\end{columns}
\end{frame}

\section{6. Сравнение
подходов}\label{ux441ux440ux430ux432ux43dux435ux43dux438ux435-ux43fux43eux434ux445ux43eux434ux43eux432}

\begin{frame}{6.1 Альтернативы --- кратко}
\phantomsection\label{ux430ux43bux44cux442ux435ux440ux43dux430ux442ux438ux432ux44b-ux43aux440ux430ux442ux43aux43e}
\begin{columns}[T]
\begin{column}{1\linewidth}
\begin{longtable}[]{@{}
  >{\raggedright\arraybackslash}p{(\linewidth - 6\tabcolsep) * \real{0.2963}}
  >{\raggedright\arraybackslash}p{(\linewidth - 6\tabcolsep) * \real{0.2222}}
  >{\raggedright\arraybackslash}p{(\linewidth - 6\tabcolsep) * \real{0.2222}}
  >{\raggedright\arraybackslash}p{(\linewidth - 6\tabcolsep) * \real{0.2593}}@{}}
\toprule\noalign{}
\begin{minipage}[b]{\linewidth}\raggedright
Подход
\end{minipage} & \begin{minipage}[b]{\linewidth}\raggedright
Суть
\end{minipage} & \begin{minipage}[b]{\linewidth}\raggedright
Плюс
\end{minipage} & \begin{minipage}[b]{\linewidth}\raggedright
Минус
\end{minipage} \\
\midrule\noalign{}
\endhead
\textbf{Урновая (наш)} & α·Sn в числителе & Прозрачно, параметр памяти &
Параметры феноменологичны \\
\textbf{Пороговая} & Фиксация после порога θ & Просто & Порог надо знать
заранее \\
\textbf{RL (кумулятивная)} & P \textasciitilde{} exp(β·Q) & Гибкая &
Сложный анализ \\
\bottomrule\noalign{}
\end{longtable}

\textbf{Наш подход --- оптимальный баланс простоты и выразительности}
\end{column}
\end{columns}
\end{frame}

\section{7. Ограничения}\label{ux43eux433ux440ux430ux43dux438ux447ux435ux43dux438ux44f}

\begin{frame}{7.1 Что модель НЕ умеет}
\phantomsection\label{ux447ux442ux43e-ux43cux43eux434ux435ux43bux44c-ux43dux435-ux443ux43cux435ux435ux442}
\begin{columns}[T]
\begin{column}{1\linewidth}
\begin{itemize}
\tightlist
\item
  \textbf{Дискретное время} → не для непрерывных процессов
\item
  \textbf{Феноменологические параметры} → не выводятся из первых
  принципов
\item
  \textbf{Нет антагонизма} → нет активного подавления конкурента
\item
  \textbf{Нужна полная история} → в реальности данные часто неполные
\item
  \textbf{Требуются длинные данные} → при n \textless{} 50 выводы
  ненадёжны
\end{itemize}

\textbf{Но:} в своих границах модель работает отлично!
\end{column}
\end{columns}
\end{frame}

\section{8. Выводы}\label{ux432ux44bux432ux43eux434ux44b}

\begin{frame}{8.1 Выводы}
\phantomsection\label{ux432ux44bux432ux43eux434ux44b-1}
\begin{columns}[T]
\begin{column}{1\linewidth}
\begin{enumerate}
\item
  \textbf{Path dependence ≠ марковский процесс}
\item
  \textbf{Модифицированная урна с α·Sn} --- простая и мощная
  формализация
\item
  \textbf{α = 0} → классика (нет памяти) \textbf{α \textgreater{} 0
  малое} → мягкая память \textbf{α большое} → жёсткая фиксация (QWERTY)
\item
  \textbf{Реальное применение:} фильтрация сетей, экономика, социология
\item
  \textbf{Важно:} помнить об ограничениях модели
\end{enumerate}
\end{column}
\end{columns}
\end{frame}

\begin{frame}{8.2 Итоговая мысль}
\phantomsection\label{ux438ux442ux43eux433ux43eux432ux430ux44f-ux43cux44bux441ux43bux44c}
\begin{columns}[T]
\begin{column}{1\linewidth}
\begin{quote}
Классические эргодические методы усредняют по траекториям и могут
ошибаться. Среднее не соответствует ни одной реальной траектории.
\end{quote}

\textbf{Понимание path dependence --- ключ к корректному моделированию
сложных систем.}
\end{column}
\end{columns}
\end{frame}




\end{document}
